\begin{ledgroupsized}[r]{120mm}
\footnotesize
\pstart
\noindent\textbf{Überlieferung:}
\pend
\end{ledgroupsized}
\begin{ledgroupsized}[r]{114mm}
\footnotesize
\pstart
\parindent -6mm
\makebox[6mm][l]{\textit{L}}%
Konzept:
LH~XXXV~10, 14~Bl.~9\textendash10. 
Ein Bogen~2\textsuperscript{o};
\href{\Scan{LH_35_10_14_009wz}}%
{Wasserzeichen} im Falz (Papier aus Wien);
geringfügiger Textverlust durch Abrieb auf Bl.~9~r\textsuperscript{o} infolge erneuter Faltung entlang der Blattmitte.
Zweieinhalb Seiten auf Bl.~9~r\textsuperscript{o} bis Bl.~10~r\textsuperscript{o};
Bl.~10~v\textsuperscript{o} leer.
Am oberen Rand von Bl.~9~r\textsuperscript{o}, in Rot, der Vermerk von unklarer Hand \textls{\textit{IX}} (vgl. die Träger der benachbarten Signaturen LH~XXXV~10, 14 Bl.~2 bis 13, u.\,a. 
N.\;\ref{60117},
\ref{60118},
\ref{60115},
\ref{60112},
\ref{60113},
\ref{60114},
\ref{60110},
\ref{41824}
u.\;\ref{60111}
in diesem Band).
\pend
\end{ledgroupsized}
\vspace{5mm}
\begin{ledgroup}
\footnotesize
\pstart
\noindent
\textbf{Datierungsgründe:}
Das Wasserzeichen Wiener Herkunft ist für Leibnizens Aufenthalte in Wien\indexO{Wien} auf dem Weg nach Italien (Mai 1688 bis Februar 1689) und auf der Rückkehr (30 April bis Mitte Mai 1690) belegt.
Folgende zwei Umstände erlauben, die Entstehung von N.\;\ref{60116} auf letzteren dieser Zeiträume zu verorten.
\pend
\pstart
Erstens 
nimmt Leibniz auf \refnote{FnC:35_10_14_009-010_Varignon} ausdrücklichen Bezug auf 
\indexN{\textls{Varignon}, Pierre de 1654\textendash1722}Pierre Varignons 
Analyse des dort behandelten Falls.
Diese Lösung konnte Leibniz im \cite{02028}\textit{Projet d'une nouvelle mechanique}, Paris 1687 (S.\;40\,f.), lesen, den er in Florenz oder Modena (Dezember 1689 bis 2. Februar 1690) exzerpiert hat: siehe N.\;\ref{60365} in diesem Band. 
Zwar war bereits in den \cite{01023}\textit{Acta Eruditorum} von August 1688 (S\;411--414) eine \cite{02166}Rezension erschienen, die auch auf die einschlägige Stelle  einging (S.\;413), und die Leibniz im Laufe seines ersten Wiener Aufenthaltes womöglich lesen konnte.
Allerdings zeigt
die Genauigkeit, mit der er Varignons \textit{fig.\;30} sowie seine eigentümlichen Bezeichnungen \textit{pondus EO} und \textit{superficies GH} übernimmt,
dass ihm das Buch selbst vorlag oder wenigstens seine ausführlichen Auszüge (\refpassage{37_05_217-218_8b}{37_05_217-218_8b} von N.\;\ref{60365} mit den Diagrammen \bracks{\textit{Fig.\;9}} und \bracks{\textit{Fig.\;10}}); genannte Einzelheiten hätte er der Rezension von 1688 nicht entnehmen können.
Folglich setzt N.\;\ref{60116} die Lektüre von Varignons Werk voraus.
\pend
\pstart
Zweitens lässt sich feststellen, dass Leibniz das Diagramm \bracks{\textit{Fig.\;7}} auf S.\;\pageref{Fig:35_10_14_009-010_Fig-7} 
in Anlehnung an die \bracks{\textit{Fig.\;4}} seines Konzepts \textit{Solutio problematis R.\,P.\;Vannii S.\,J.}
angefertigt hat (siehe den Zeugen \textit{L\textsuperscript{2}} von N.\;\ref{41823}, S.\;\pageref{Fig:35_10_13_010-011_Fig-4}).
Er hat zunächst im vorliegenden Konzept die Figur von N.\;\ref{41823} nachgezeichnet, auch die Buchstaben zur Bezeichnung der Punkte fast ausnahmslos übernommen; diese hat er aber durch Zahlen ersetzt und im einschlägigen Abschnitt von N.\;\ref{60116} die Punkte durch entsprechende Zahlen bezeichnet.
Daher muss N.\;\ref{60116} nach N.\;\ref{41823} entstanden sein,
zumindest nach dessen Kernteil (\refpassage{35_10_13_010-011_16a}{35_10_13_010-011_16a}), den Leibniz auf dem Rückweg nach Deutschland, vor oder spätestens während seines Aufenthaltes in Wien verfasst hatte, und den er noch im Mai 1690 zu veröffentlichen beabsichtigte (siehe die Einführung zu N.\;\ref{intro_vanni-solutio-problematis}).
\pend 
\end{ledgroup}
\clearpage
\selectlanguage{latin}
\vspace{8mm}
\pstart
\noindent
\href{\Scan{LH_35_10_14_009r}}%
{\lbrack9~r\textsuperscript{o}\rbrack} 
\pend\nobreak
\begin{edfigure}[0em][.5em]
\includegraphics[scale=1]{%
gesamttex/edit_VIII-4/images/LH_35_10_14_009-010_d1_009r.pdf}
\edcaption{%
\edtext{\lbrack\textit{Fig.~1, tlw. Blindzeichnung}\rbrack}{%
\lemma{\lbrack\textit{Fig.~1}\rbrack}%
\Cfootnote[nonum]{Der Punkt \bracks{(M)} heißt bei Leibniz \textit{M}; diese Bezeichnung ändert Hrsg. gemäß den Angaben im Text. Ein gestrichener Entwurf zum Diagramm wird nicht wiedergegeben.}}}
\end{edfigure}
\vspace{0.5em}		%
\pstart
\noindent
\lbrack\textit{Nachfolgend kleingedruckter Text gestrichen:}\rbrack
\pend
\nobreak
\vspace{0.5em}
\footnotesize	%
\pstart 
\noindent
\edtext{Sit onus \textit{GCD} cujus}{%
\lemma{Sit}%
\Bfootnote{%
\textit{(1)}~pondus
\textit{(2)}~onus \bars{\textit{GCD} \textit{erg.}} cujus~\textit{L}}}
centrum gravitatis\indexS{centrum gravitatis} \textit{G}; superficies \textit{VDS} qua tangit planum sustentans\indexS{planum sustentans} \textit{AD} punctum affixae chordae 
\edtext{\textit{CBE}, tractae}{%
\lemma{\textit{CBE},}%
\Bfootnote{%
\textit{(1)}~trahentis
\textit{(2)}~tractae~\textit{L}}}
in linea \textit{BC}, a pondere \textit{E} appenso,\indexS{pondus appensum} et perpendiculariter descendente.
\edtext{Quaeritur magnitudo}{%
\lemma{Quaeritur}%
\Bfootnote{%
\textit{(1)}~\bars{situs \textit{streicht Hrsg.}} talis
\textit{(2)}~pondus
\textit{(3)}~magnitudo~\textit{L}}}
ponderis \textit{E}, ut sustineat onus 
\edtext{\textit{CGD} et}{%
\lemma{\textit{CGD}}%
\Bfootnote{%
\textit{(1)}~nec tamen
\textit{(2)}~etiam plus
\textit{(a)}~vel
\textit{(b)}~vel minus
\textit{(c)}~et~\textit{L}}}
ab ipso vicissim sustineatur.\indexS{pondera vicissim sustinentia}
\pend
\normalsize		%
\vspace{0.5em}		%
\pstart
\noindent
\edtext{Sit onus \textit{GCSD\greek{d}V}}{%
\lemma{Sit}%
\Bfootnote{%
\textit{(1)}~\textit{G}
\textit{(2)}~onus
\textit{(a)}~\textit{GCD}
\textit{(b)}~\textit{GCSD\greek{d}V}~\textit{L}}}
cujus centrum gravitatis\indexS{centrum gravitatis} \textit{G}, superficies \textit{VDS} qua tangit planum sustentans\indexS{planum sustentans} \textit{AD} in puncto \textit{D}; retineatur praeterea ab affixa chorda \textit{CBE}, tracta in linea \textit{CB} super trochlea\indexS{trochlea} \textit{B}, a pondere\indexS{pondus chordam super trochlea trahens} \textit{E}, perpendiculariter deorsum nitente in recta \textit{BE}, cui (si opus productae\bracks{)} 
\edtext{planum sustentans\indexS{planum sustentans}}{%
\lemma{planum}%
\Bfootnote{%
\textit{(1)}~occurrat in
\textit{(2)}~sustentans~\textit{L}}}
occurrat in \textit{A}. Ponamus jam volvi onus\indexS{onus super plano volvens} super plano, 
\edtext{et ita}{%
\lemma{et}%
\Bfootnote{%
\textit{(1)}~ita %
\textit{(2)}~p %
\textit{(3)}~cum %
\textit{(4)}~ita~\textit{L}}}
transferri ex 
\edtext{situ \textit{GCSD\greek{d}V}}{%
\lemma{situ}%
\Bfootnote{%
\textit{(1)}~\textit{GCD\greek{d}}, %
\textit{(2)}~\textit{GCD\greek{d}V},
\textit{(2)}~\textit{GCSD\greek{d}V}~\textit{L}}}
in situm \textit{(G)(C)(S)(D)(\greek{d})(V)} planumque jam tangere in 
\edtext{puncto \textit{(\greek{d})} atque ita puncto affixionis\indexS{punctum affixionis} \textit{C} translato}{%
\lemma{puncto}%
\Bfootnote{%
\textit{(1)}~\textit{(D)} %
\textit{(2)}~\textit{(\greek{d})} atque ita
\textit{(a)}~\textit{C} translato in \textit{(C)}, pondus \textit{E} %
\textit{(b)}~puncto affixionis \bars{\textit{C} \textit{erg.}} translato~\textit{L}}}
in \textit{(C)} et centro gravitatis\indexS{centrum gravitatis} \textit{G}, descendente in \textit{(G)} pondus \textit{E}
\edtext{ascendisse in \textit{(E)} debent compensari}{%
\lemma{ascendisse}%
\Bfootnote{%
\bars{in \textit{(E)} \textit{erg.}} %
\textit{(1)}~debet, %
\textit{(2)}~\bars{debent, \textit{ändert Hrsg.}} compensari~\textit{L}}}
inter se ascensus\indexS{ascensus ponderis} descensusque,\indexS{descensus ponderis} ita ut in neutram partem sit 
lucrum; hoc est ut altitudo\indexS{altitudo descensus}
descensus, \textit{Q(Q)} ducta in pondus \textit{G}, aequetur altitudini\indexS{altitudo ascensus} \textit{E(E)} ductae in pondus\indexS{altitudo ducta in pondus} \textit{E}; 
adeoque quod eodem redit, centrum gravitatis commune\indexS{centrum gravitatis commune} 
\textit{K} neque ascenderit neque descenderit, adeoque nec punctum \textit{R} ei respondens in recta 
\edtext{perpendiculari \textit{BA}. Evanescente proinde}{%
\lemma{perpendiculari \textit{BA}.}%
\Bfootnote{%
\textit{(1)}~Quod %
\textit{(2)}~Evanescente %
\textit{(a)}~differentia inter %
\textit{(b)}~proinde~\textit{L}}}
\textit{R(R)} seu differentia inter \textit{BR} 
\edtext{et \textit{B(R)},}{%
\lemma{et}%
\Bfootnote{%
\textit{(1)}~inter \textit{TB} %
\textit{(2)}~\textit{B(R)},~\textit{L}}}
quod idem est, ac si diceremus 
\edtext{\textit{BR} vel \textit{\greek{W}R} quaeri maximam vel minimam. Itaque}{%
\lemma{}%
\Bfootnote{%
\textit{BR} \bars{vel \textit{\greek{W}R}~\textit{erg.}} quaeri maximam \bars{vel minimam~\textit{erg.}}. Itaque~\textit{L}}}
perinde est ac si calculum instituamus quem hactenus exposuimus; efficiendo, 
\edtext{ut fiat}{%
\lemma{ut}%
\Bfootnote{%
\textit{(1)}~fiant
\textit{(2)}~fiat~\textit{L}}}
\textit{E(E)} in pondus \textit{E} aequalis ipsi \textit{Q(Q)} ductae in pondus \textit{G}; an vero quod longe est 
\edtext{compendiosius; si}{%
\lemma{compendiosius;}%
\Bfootnote{%
\textit{(1)}~invento valore ipsius \textit{BR} %
\textit{(2)}~si~\textit{L}}}
ex assumta \textit{E} quacunque, et simplici situ adhibito ipsius oneris\indexS{situs oneris}
\edtext{\textit{GCSD\greek{d}V}, valorem ipsius \textit{BR}}{%
\lemma{\textit{GCSD\greek{d}V},}%
\Bfootnote{%
\textit{(1)}~calculo investige %
\textit{(2)}~calculo exper %
\textit{(3)}~valorem ipsius \textit{BR}~\textit{L}}}
analytice inventum\indexS{valor analytice inventus} faciamus
\edtext{o\bracks{m}nium}{%
\lemma{}%
\Bfootnote{%
ommnium \textit{L ändert Hrsg.}}} 
possibilium Maximum,\indexS{valor omnium possibilium maximus}
\edtext{adhibitis meis artificiis,\indexS{artificia mea de maximo et minimo} de maximo et minimo,}{%
\lemma{adhibitis \bracks{...} minimo}%
\Cfootnote{Siehe \textsc{G.\,W.~Leibniz}, \cite{02077}\textit{Nova methodus pro maximis et minimis}, \cite{01023}\textit{AE} Oktober 1684, S.\;467--473.}}
quae nullas irrationales morantur, et mire calculum contrahunt.
Ubi tamen inprimis notandum est, si \textit{G} immineat ipsi \textit{D} nullam fieri volutionem,\indexS{volutio}
sed simplicem defluxionem\indexS{defluxio} manente eodem puncto tangente \textit{D} 
\edtext{in quo casu calculum jam sum executus.}{%
\lemma{in \bracks{...} executus}%
\Cfootnote{%
Siehe bspw. \textit{Solutio problematis R.\,P.\;Vannii S.\,J.}, N.\;\ref{41823}, bes. \refpassage{35_10_13_010-011_9a}{35_10_13_010-011_9b} im Zeugen \textit{L\textsuperscript{1}}.}}
Video praeterea rem his propositis non satis fore determinatam; prout enim initium assumitur volutionis,\indexS{volutio} aliud atque aliud prodibit
\edtext{maximum; ut si}{%
\lemma{maximum;}%
\Bfootnote{%
\textit{(1)}~imo assumto quolibet \bars{puncto \textit{D} \textit{streicht Hrsg.}} inveniri potest situs %
\textit{(a)}~ejus optimus %
\textit{(b)}~in plano talis, ut %
\textit{(2)}~ut si~\textit{L}}}
\textit{S} primum cadat in\;\textit{A}, volvetur curva, descendensque inveniet tandem situm aequilibrii\indexS{situs aequilibrii} attrahens pondus\;\textit{E}.
\pend
\pstart
Itaque omissa consideratione curvae, optima methodus haec esse videtur ut rem staticam\indexS{res statica seu res aequilibrii} seu aequilibrii,\indexS{aequilibrium} et quid sustentet, adeoque quid moveat\bracks{,} 
reducamus primum ad apices\indexS{apex grave sustentans} quosdam perpendiculari nisu\indexS{nisus perpendicularis} gravia sustentantes, ita ut nec superficiei gravis,\indexS{superficies gravis} nec superficiei sustentantis, nec lineae 
\edtext{adeo trahentis}{%
\lemma{adeo}%
\Bfootnote{%
\textit{(1)}~nisus %
\textit{(2)}~illi %
\textit{(3)}~trahentis~\textit{L}}}
vel retinentis habeatur 
\edlabel{35_10_14_009-010_1a}%
\edtext{}{%
\xxref{35_10_14_009-010_1a}{35_10_14_009-010_1b}%
\lemma{ratio.}%
\Bfootnote{%
\textit{(1)}~Sit corpus grave, cujus centrum gravitatis \textit{G} %
\textit{(a)}~, sustentatum ab %
\textit{(b)}~. \bars{Hoc \textit{streicht Hrsg.}} corpus sustentetur ab apicibus \bars{seu punctis \textit{streicht Hrsg.}} \textit{A} et \textit{B}, 
\textit{(2)}~Sit gravis cen %
\textit{(3)}~Sit~\textit{L}}}%
gravis centrum \textit{G}, ratio.
\pend
\begin{edfigure}%
\includegraphics[scale=1]{%
gesamttex/edit_VIII-4/images/LH_35_10_14_009-010_d2_009r.pdf}
\edcaption[.2em]{\edtext{\lbrack\textit{Fig.~2}\rbrack}{%
\lemma{\lbrack\textit{Fig.~2}\rbrack}%
\Cfootnote[nonum]{%
Drei gestrichene Entwürfe zum Diagramm werden nicht wiedergegeben.}}}
\end{edfigure}
\pstart
Sit\edlabel{35_10_14_009-010_1b}
gravis centrum \textit{G}, id innitatur puncto \textit{D}, sed pro parte tantum, ut rotam scilicet circa \textit{D}, adeoque descendere possit, quaeritur quantum sit pondus retinens\indexS{pondus grave retinens} \textit{E}.
\pend
\pstart
Cum pondus\indexS{pondus planis incumbens} 
\edtext{cujus centrum gravitatis\indexS{centrum gravitatis} est}{%
\lemma{}%
\Bfootnote{%
cujus centrum gravitatis est \textit{erg.~L}}}
\textit{G} incumbit planis \textit{AB}, \textit{EF}, punctis \textit{C} et \textit{D}, perinde est ac si punctum seu centrum \textit{G} impellatur rectis 
\textit{CG}, \textit{DG} tanquam contis;\indexS{contus}\indexS{conti centrum gravitatis impellentes} sed hi ipsi conti\indexS{contus} impellantur 
\edtext{nisibus\indexS{nisus} \textit{LC}, \textit{MD}.}{%
\lemma{nisibus \textit{LC},}%
\Bfootnote{%
\textit{(1)}~\textit{DM} %
\textit{(2)}~\textit{MD}.~\textit{L}}}
Nimirum grave \textit{G} tendens conatu\indexS{conatus} \textit{GK} perpendiculari tendit simul conatibus \textit{GD}, \textit{GN}\bracks{,} tendens autem conatu \textit{GD} vel \textit{DR}, premit planum \textit{FE} conatu \textit{DM}, tendensque conatu \textit{GN} vel \textit{CS} premit planum \textit{AB} conatu \textit{CL}. %
\pend
\begin{edfigure}[.2em][1em]
\includegraphics[scale=1]{%
gesamttex/edit_VIII-4/images/LH_35_10_14_009-010_d3_009r.pdf}
\edcaption{\edtext{\lbrack\textit{Fig.~3}\rbrack}{%
\lemma{\lbrack\textit{Fig.~3}\rbrack}%
\Cfootnote[nonum]{%
Drei gestrichene Entwürfe zum Diagramm werden nicht wiedergegeben.}}}
\end{edfigure}
\pstart
Videtur hic nasci improba quaedam difficultas; nitens 
\edtext{in \textit{EF} nisu \textit{GD}, vel \textit{DR}, afficit}{%
\lemma{in \textit{EF}}%
\Bfootnote{%
\textit{(1)}~impetu \textit{GD} %
\textit{(2)}~nisu \textit{GD}, 
\bars{\textit{(1)}~vel \textit{DR} %
\textit{(2)}~vel \textit{DR}, \textit{erg.}}
afficit~\textit{L}}}
ipsum tantum nisu \textit{DM}. Ergo videtur resistentia\indexS{resistentia} opus esse ut \textit{DM}; 
\edtext{adeoque sufficiet}{%
\lemma{adeoque}%
\Bfootnote{%
\textit{(1)}~pondere %
\textit{(2)}~sufficiet~\textit{L}}}
pondus ut \textit{D}\textit{M}. Verum vicissim pondus agens directione \textit{MD}; dat ipsi \textit{G}, in directione \textit{DG} tantum conatum\indexS{conatus} ut \textit{VD}, non vero ut \textit{RD}, atque ita ad imprimendum conatum ut \textit{RD}, opus videtur pondere ut \textit{TD}. Quomodo igitur dicemus premi planum \textit{EF}, ponderene 
\edtext{ut \textit{DM}}{%
\lemma{}%
\Bfootnote{%
ut \bars{ut \textit{streicht Hrsg.}} \textit{DM}~\textit{L}}}
an ut 
\edtext{\textit{DT}, pondus vero resistens debere esse}{%
\lemma{\textit{DT},}%
\Bfootnote{%
\textit{(1)}~et pon
\textit{(2)}~pondus vero resistens
\textit{(a)}~premi nisu et %
\textit{(b)}~debere esse~\textit{L}}}
\ledpb%
ut \textit{TD}, an solum ut \textit{MD}. An dicemus aliam esse vim quam excipit planum,\indexS{vis quam excipit planum} 
quae sit ut \textit{DM}, aliam esse vim quam debeat habere pondus quae sit ut \textit{TD}. Cum tamen rationis 
\edtext{sit, utramque vim}{%
\lemma{sit,}%
\Bfootnote{%
\textit{(1)}~utraque vis %
\textit{(2)}~utramque vim~\textit{L}}}
quam patitur planum,\indexS{vis quam patitur planum} et qua sustineri potest\indexS{vis qua planum sustineri potest} contra actionem,\indexS{actio} esse eandem.
\pend
\pstart
Concipi etiam potest, pondus \textit{G} circa centrum \textit{C} agere vectem\indexS{vectis} \textit{CGD}
eodemque propellere punctum \textit{D}, indeque comparandus est descensus 
\edtext{puncti \textit{G} seu progressus}{%
\lemma{puncti \textit{G}}%
\Bfootnote{%
\textit{(1)}~cum propulsione %
\textit{(2)}~seu progressus~\textit{L}}}
in parallela ad \textit{GK}, cum progressu
\edtext{puncti \textit{D} in}{%
\lemma{puncti}%
\Bfootnote{%
\textit{(1)}~cum %
\textit{(2)}~\textit{D} in~\textit{L}}}
parallela ad \textit{D}\textit{M}. Est autem descensus ipsius \textit{G} ad descensum ipsius \textit{D} (in parallelis ad \textit{GK}) ut \textit{CN} ad \textit{CD}.
\pend
\pstart
\href{\Scan{LH_35_10_14_009v}}%
{\lbrack9~v\textsuperscript{o}\rbrack}\ 
Videndum quomodo corpus cujus centrum gravitatis\indexS{centrum gravitatis} \textit{G}, innitens plano\indexS{planum inclinatum} \textit{AD}, sed cum libertate descendendi\indexS{libertas descendendi} moveatur 
\edtext{in plano, et}{%
\lemma{in plano,}%
\Bfootnote{%
\textit{(1)}~et simul
\textit{(2)}~et~\textit{L}}}
an ac quatenus simul volvatur\indexS{corpus volvens} et 
\edlabel{35_10_14_009-010_2a}%
\edtext{}{%
\xxref{35_10_14_009-010_2a}{35_10_14_009-010_2b}%
\lemma{\textls{defluat}. \bracks{/}}%
\Bfootnote{%
\textit{(1)}~Pondus \textit{EO} %
\textit{(2)}~Si Varignonio credimus pondus \textit{EO},~\textit{L}}}%
\textls{defluat}.
\pend
\begin{edfigure}%
\hfill
\includegraphics[scale=1]{%
gesamttex/edit_VIII-4/images/LH_35_10_14_009-010_d4_009v.pdf}
\hfill%
\includegraphics[scale=1]{%
gesamttex/edit_VIII-4/images/LH_35_10_14_009-010_d5_009v.pdf}
\hfill
\edcaption{%
\hspace*{10mm}%
\lbrack\textit{Fig.~4}\rbrack 
\hspace*{30mm}%
\edtext{\lbrack\textit{Fig.~5, tlw. Blindzeichnung}\rbrack%
}{%
\lemma{\lbrack\textit{Fig.~5}\rbrack}%
\Cfootnote[nonum]{%
Vorlage ist die "fig.\;30"
bei \indexN{\textls{Varignon}, Pierre de 1654\textendash1722}\textsc{P.~Varignon},
\cite{02028}\textit{Projet d'une nouvelle mechanique}, Paris 1687, Tafel 4.}}}
\end{edfigure}
\pstart
\edtext{Si \indexN{\textls{Varignon}, Pierre de 1654\textendash1722}Varignonio credimus resistentia superficiei\indexS{resistentia superficiei} \textit{GH}, et potentia\indexS{potentia} \textit{R}, 
sunt invicem, ut lineae \textit{AC}, \textit{AD}, \textit{AB}.}{%
\applabel{FnC:35_10_14_009-010_Varignon}%
\lemma{Si \bracks{...} \textit{AC}, \textit{AD}, \textit{AB}}%
\Cfootnote{Es handelt sich um \indexN{\textls{Varignon}, Pierre de 1654\textendash1722}Varignons \textit{Proposition fondamentale des poids soutenus sur quelque espèce de surfaces que ce soit}, \cite{02028}a.\,a.\,O., S.\;40\,f.}}
pondus\indexS{pondus} \textit{EO},\edlabel{35_10_14_009-010_2b}
Posito \textit{A} esse 
punctum concursus directionis ponderis\indexS{directio ponderis} et potentiarum\indexS{concursus directionis potentiarum et ponderis}
praeterea necessario debet esse \textit{AD} perpendicularis ad 
\edtext{\bracks{\textit{GH}}}{\lemma{}\Bfootnote{\textit{G} \textit{L ändert Hrsg.}}}
in \textit{O}, alioqui corpus volvetur super superficie.\indexS{corpus volvens super superficie} 
Haec itaque perinde ipsi 
\edtext{est, in}{%
\lemma{est,}%
\Bfootnote{%
\textit{(1)}~utrum centrum grav %
\textit{(2)}~in~\textit{L}}}
quo puncto rectae \textit{AC} centrum 
\edtext{gravitatis\indexS{centrum gravitatis} corporis}{%
\lemma{gravitatis}%
\Bfootnote{%
\textit{(1)}~sit %
\textit{(2)}~corporis~\textit{L}}}
sit situm. Et quidem videtur rationi consentaneum nam perinde est, si esset linea 
\ledpb%
rigida pondere carens \textit{AO}, et punctum \textit{A} traheretur appenso pondere\indexS{pondus appensum} \textit{F}, (ubi nihil refert quanta sit longitudo ipsius
\edtext{\textit{AF}) et}{%
\lemma{\textit{AF})}%
\Bfootnote{%
\textit{(1)}~punctum autem %
\textit{(2)}~et~\textit{L}}}
potentia\indexS{potentia} \textit{R}.
\pend
\begin{edfigure}%
\hfill
\includegraphics[scale=1]{%
gesamttex/edit_VIII-4/images/LH_35_10_14_009-010_d6_009v.pdf}
\hfill%
\includegraphics[scale=1]{%
gesamttex/edit_VIII-4/images/LH_35_10_14_009-010_d7_009v.pdf}
\hfill
\edcaption{%
\pagelabel{Fig:35_10_14_009-010_Fig-7}%
\edtext{\lbrack\textit{Fig.~6}\rbrack}{%
\lemma{\lbrack\textit{Fig.~6}\rbrack}%
\Cfootnote[nonum]{%
Leibniz hat die Bezeichnung "\textit{22}" an zwei verschiedene Punkte vergeben.}} 
\hspace*{52mm}%
\edtext{\lbrack\textit{Fig.~7}\rbrack%
\hspace*{8mm}}{%
\lemma{\lbrack\textit{Fig.~7}\rbrack}%
\Cfootnote[nonum]{%
Zur Bezeichnung der Punkte hat Leibniz zunächst Buchstaben verwendet, die fast ausnahmslos aus dem Diagramm \bracks{\textit{Fig.\;4}} von N.\;\ref{41823} (S.\;\pageref{Fig:35_10_13_010-011_Fig-4}) übernommen waren, und erst später die Buchstaben durch Zahlen ersetzt (z.\,B. \textit{4}, \textit{5}, \textit{6}, \textit{7} statt \textit{F}, \textit{H}, \textit{E}, \textit{Q}).}}}
\end{edfigure}
\pstart
Sed si pondus simul pluribus planis incumbat,\indexS{pondus pluribus planis incumbens}
ut in \textit{N} et \textit{O}, major est 
\edtext{difficultas, cum}{%
\lemma{difficultas,}%
\Bfootnote{%
\textit{(1)}~quidem si %
\textit{(2)}~cum~\textit{L}}}
rectae \textit{N.1}, \textit{F.2}, \textit{O.3} non concurrant in uno puncto, sed in tribus \textit{4},\;\textit{5}\;\textit{6}; quae si coinciderent cessat difficultas; sed aliis casibus manet. 
Videtur sic procedi posse: ob directionem resistentiae\indexS{directio resistentiae} in \textit{N} quae est \textit{N.4.7}, et directionem resistentiae\indexS{directio resistentiae} in \textit{O}, quae est
\edtext{\textit{O.6.8}, a quibus}{%
\lemma{\textit{O.6.8},}%
\Bfootnote{%
\textit{(1)}~perinde est %
\textit{(2)}~a quibus~\textit{L}}}
linea directionis\indexS{linea directionis ponderis} ipsius 
\edtext{ponderis nempe \textit{6.4.F},}{%
\lemma{ponderis}%
\Bfootnote{%
\textit{(1)}~secatur in
\textit{(2)}~nempe \textit{6.4.F},~\textit{L}}}
secatur in \textit{6} et \textit{4}, perinde est
ac si adesset linea rigida \textit{6.4.F}, pondere carens, nisi quod totum pondus collectum seu appensum\indexS{pondus appensum} in \textit{F}, et haec linea rigida 
\edtext{trahatur chordis}{%
\lemma{trahatur}%
\Bfootnote{%
\textit{(1)}~in \textit{8} %
\textit{(2)}~chordis~\textit{L}}}
\textit{8.6}, et \textit{7.4}, quaeritur quae debeat esse proportio ponderum,\indexS{proportio ponderum} ut \textit{6.4} 
\edtext{maneat horizonti perpendicularis. Pro}{%
\lemma{maneat}%
\Bfootnote{%
\textit{(1)}~linea %
\textit{(2)}~recta per %
\textit{(3)}~horizonti perpendicularis. %
\textit{(a)}~Tollamus %
\textit{(b)}~Pro~\textit{L}}}
uno pondere\indexS{pondus} concipiamus tantum chordam seu potius lineam rigidam \textit{8.6}, in \textit{7} firmiter annexam et circa \textit{7} rotatilem. Pondere 
\edtext{jam altero}{%
\lemma{jam}%
\Bfootnote{%
\textit{(1)}~assumto %
\textit{(2)}~altero~\textit{L}}}
\textit{9} assumto quocunque quaeritur situs ipsius \textit{8.6}, adeoque et ipsius \textit{6.4.F} (semper verticaliter ponendi) talis,
ut descensus centri gravitatis\indexS{centrum gravitatis}\indexS{descensus centri gravitatis} \textit{9} et \textit{F}, 
sit omnium possibilium maximus,\indexS{descensus centri gravitatis omnium possibilium maximum} id centrum sit \textit{K}, debet \textit{7.K} esse maxima. 
Datur \textit{8.6} licet pro arbitrio assumenda, datur et \textit{8.10}, et \textit{7.10}, item \textit{6.4}, et \textit{6.F}. 
\pend
\pstart
Assumatur \textit{7.11} pro arbitrio, datur et \textit{11.10}, quae est $\textit{7.10}-\textit{7.11}$ cujus quadratum si subtrahas ab \textit{8.6} data, habetur \textit{8.12}, quam subtrahendo ab \textit{8.10}, restat \textit{10.12} 
\edlabel{35_10_14_009-010_4a}%
\edtext{}{%
\xxref{35_10_14_009-010_4a}{35_10_14_009-010_4b}%
\lemma{}%
\Afootnote{%
\textit{Am Rand}: Hic calculus supponit rectam \textit{6.4.F} vi aliqua inter mutandum horizonti\textsuperscript{\bracks{a}} perpendicularem servari et ita quaeri optimum situm ita non opus supponere semper \textit{6.4.F} maneret perpendicularem.
\\[1.5ex]
{\footnotesize
\textsuperscript{\lbrack a\rbrack}horizonti %
\textit{(1)}~paralle %
\textit{(2)}~recta %
\textit{(3)}~perpendicularem~\textit{L}}}}%
\edtext{vel \textit{6.11} vel \textit{7.13},}{%
\lemma{vel \textit{6.11}}%
\Bfootnote{%
\textit{(1)}~v %
\textit{(2)}~porro
\textit{(3)}~vel \textit{7.13},~\textit{L}}}
cujus quadrato si addas qu. \textit{4.13} (quae fit a \textit{6.4} data subtrahendo \textit{10.11} seu \textit{6.13}) restat quadr. de \textit{4.7}. Jam si \textit{4.7} subtrahas a $\textit{4.7}+\textit{7.9}$ data
\edtext{restat \textit{7.9}. Jam $\textit{13}.F-\textit{7.9}=\textit{9.14}$ seu $\textit{6}.F-\textit{7.10}$}{%
\lemma{restat \textit{7.9}.}%
\Bfootnote{%
\textit{(1)}~cui addas \textit{9.K}. quae se habet %
\textit{(2)}~Jam %
\textit{(a)}~$\textit{7.9}-\textit{13}.F$ seu \textit{7.9} %
\textit{(b)}~$\textit{13}.F-\textit{7.9}$~\textit{L}}}
$\!-\,\textit{7.9}=\textit{9.14}$,
jam \textit{9.K} est ad \textit{9.14} 
\edtext{ut pondus}{%
\lemma{ut}%
\Bfootnote{%
\textit{(1)}~\textit{F} ad
\textit{(2)}~pondus~\textit{L}}}
\textit{F} ad pondus $\textit{9}+F$, addendo ergo \textit{9.K} ad \textit{7.9}. habetur \textit{7.K} quaesita ad minimam
\edlabel{35_10_14_009-010_3a}%
\edtext{}{%
\xxref{35_10_14_009-010_3a}{35_10_14_009-010_3b}%
\lemma{determinanda.}%
\Bfootnote{%
\textit{(1)}~\textit{7.11} sit %
\textit{(a)}~\textit{y} %
\textit{(b)}~\textit{7}. \textit{8.6} sit %
\textit{(aa)}~\textit{8} %
\textit{(bb)}~\textit{6} %
\textit{(cc)}~\textit{8}; \textit{6.F} sit \textit{6}. \textit{6.4} sit \textit{4}. %
\textit{(aaa)}~$\textit{4.7}+\textit{7}$ %
\textit{(bbb)}~\textit{8.10} sit \textit{10}. \textit{7.10} sit \textit{12}. $\sqrt{-\overline{\textit{12}-\textit{7}}^2+\textit{8}^2}=\textit{8}.\textit{12}$ %
\textit{(2)}~\textit{7.11} sit \textit{y}~\textit{L}}}%
determinanda.
\pend
\pstart
\textit{7.11} sit \textit{y}\bracks{,}\edlabel{35_10_14_009-010_3b}
\textit{8.6},\,\textit{a}\bracks{,} \textit{4.7.9}\,\textit{e}, 
\textit{8.10}\,\textit{b}, \textit{7.10},\,\textit{c}.\bracks{,} 
\textit{6.4}\,\textit{g}, \textit{6.F}\,\textit{h}\bracks{,} pondus \textit{9},\,\textit{p},
pondus \textit{F}\,\textit{f}. %
\pend
\pstart
$\textit{11.10}=c-y$,
\edtext{fit}{\lemma{}\Bfootnote{fit \textit{erg.~L}}}
$\sqrt{aa-\overline{c-y}^2}=\textit{8.12}=m$\bracks{,}
$b-m=\textit{10.12}$ vel \textit{6.11} vel \textit{7.13}\bracks{,} $\textit{4.13}=\textit{6.4}-\textit{7.11}=g-y$\bracks{,} 
$\overline{7.13}^2+\overline{4.13}^2=\overline{4.7}^2$.
Ergo $\sqrt{\overline{b-m}^2+\overline{g-y}^2}=\overline{4.7}=n$, $e-n=\textit{7.9}$. 
\pend
\pstart
Porro $\textit{13}.F=\textit{6}.F-7.\textit{10}=h-y=\textit{7.14}$, unde auferendo \textit{7.9}. 
\edtext{fit $h-y-e+n=\textit{9.14}$}{%
\lemma{fit $h-y-e+n=$}%
\Bfootnote{%
\textit{(1)}~\textit{K.14} %
\textit{(2)}~\textit{7.14} %
\textit{(3)}~\textit{9.14}~\textit{L}}}
cujus%
\edlabel{35_10_14_009-010_4b}
sumenda portio \textit{9.K}, quae fit multiplicando per \textit{f}, et dividendo per $f+p$. Ergo fiet $\textit{9}.K=\overline{h-y-e+n}\;f:\overline{f+p}$ quam addendo ad \textit{7.9} 
fit $e-n+\overline{\overline{h-y-e+n}\;f:\overline{f+p}}= \text{Max}$.
\pend
\begin{edfigure}%
\includegraphics[scale=1]{%
gesamttex/edit_VIII-4/images/LH_35_10_14_009-010_d8_009v.pdf}
\edcaption[.2em]{\edtext{\lbrack\textit{Fig.~8}\rbrack}{%
\lemma{\lbrack\textit{Fig.~8}\rbrack}%
\Cfootnote[nonum]{%
Ein nicht gestrichener Entwurf zum Diagramm wird nicht wiedergegeben.}}}
\end{edfigure}
\pstart
Unde fiet $d\overline{n}\;\overline{f+p}=dnf-d\overline{y}\; f$ seu $d\overline{n}\; p+d\overline{y}\; f=0$, idque praevideri poterat, sine calculo, nihil aliud enim dicit quam ascensum\indexS{ascensus ponderis} 
\edtext{ipsius puncti \textit{6},}{%
\lemma{ipsius}%
\Bfootnote{%
\textit{(1)}~ponderis \textit{6}, qui est $d\overline{n}$ %
\textit{(2)}~puncti \textit{6},~\textit{L}}}
seu ponderis \textit{f} qui est $d\overline{y}$, et descensum ponderis\indexS{descensus ponderis} \textit{p}, qui est $d\overline{n}$, vel potius qui est $-d\overline{n}$ qui in 
\edtext{quantum \textit{n} seu}{%
\lemma{quantum}%
\Bfootnote{%
\textit{(1)}~\textit{n}, crescit, in ta %
\textit{(2)}~\textit{n} seu~\textit{L}}}
\textit{4.7} crescit, in tantum \textit{7.9} decrescit et contra, esse reciproce ut pondera.
\edtext{Jam $d\overline{n}=\overline{\cancel{2}\;\overline{b-m}.\overline{-dm}+\cancel{2}\;\overline{g-y}.\overline{-dy}}:\cancel{2}n$}{%
\lemma{}%
\Bfootnote{%
Jam $d\overline{n}$ \bars{seu potius $2d\overline{n}$ \textit{gestr.}} $=\overline{\cancel{2}\;\overline{b-m}.\overline{-dm}+\cancel{2}\;\overline{g-y}.\overline{-dy}}:\cancel{2}n$~\textit{L}}}
et $d\overline{m}=-2\;\overline{c-y}.\overline{-dy}:2m=\overline{c-y}\; dy:m$. Fit $-\overline{b\cdot\overline{c-y}:m}+
\underset{c-g}
{\ovalbox{\overline{c-y}-\overline{g-y}}}
=-fn:p$.
\pend
\pstart
Jam \textit{n} 
\edtext{e\bracks{s}t}{%
\lemma{}\Bfootnote{et \textit{L ändert Hrsg.}}}
\textit{4.7}. $c-y$ est \textit{10.11}. vel \textit{6.12} et \textit{m} est \textit{8.12}. Hinc  
\edtext{$b\cdot\overline{c-y}:m$ est recta}{%
\lemma{$b\cdot\overline{c-y}:m$ est}%
\Bfootnote{%
\textit{(1)}~quantitas %
\textit{(2)}~recta~\textit{L}}}
quae est ad \textit{b} seu \textit{8.10} ut $c-y$ seu \textit{6.12} est ad \textit{m} seu ad \textit{8.12}. Sed haec recta est \textit{10.15}. quae habetur producendo \textit{86}. dum ipsi \textit{10.14} occurrat in \textit{15}, cui si addas \textit{15.16} seu \textit{g}, quae habetur ducendo \textit{4.16} parallelam ipsi \textit{8.6}. ipsi \textit{10.14} occurrentem in \textit{16}, fiet $\textit{10.15}=\overline{b\cdot\overline{c-y}:m}+g$, unde si auferas \textit{7.10} seu \textit{c}, fiet $\textit{7.16}=b.\overline{c-y}:m+g-c$. Ergo fit $\textit{7.16}=fn:p$ seu \textit{7.16} ad \textit{4.7}, ut \textit{f} ad \textit{p}. Jam \textit{8.15} secet \textit{4.7} in \textit{5}, patet esse \textit{7.16} ad \textit{7.4} ut \textit{15.16} seu \textit{6.4} ad \textit{5.4}. %
\pend
\pstart
Itaque pondus appendendum\indexS{pondus appensum} in \textit{9} est ad 
\edtext{pondus\indexS{pondus} datum}{%
\lemma{pondus}%
\Bfootnote{%
\textit{(1)}~appendend %
\textit{(2)}~datum~\textit{L}}}
integrum \textit{f}, ut \textit{4.5} ad \textit{4.6} datam.%
\pend
\pstart
Eodem ergo modo investigabimus quantitatem ponderis\indexS{pondus appensum} in \textit{18} appendendi, producatur \textit{7.4} donec ipsi \textit{8.18} occurrat in \textit{19}. et huic per \textit{6} ducatur parallela \textit{6.20}, erit pondus \textit{18} ad pondus \textit{F} ut \textit{8.6} ad \textit{8.20}, seu 
\edtext{ut \textit{6.5}}{%
\lemma{ut}%
\Bfootnote{%
\textit{(1)}~\textit{6.17} %
\textit{(2)}~\textit{6.5}~\textit{L}}}
ad \textit{6.4}. 
Itaque pulcherrima conclusione, latera 
\edtext{trianguli \textit{6.4.5}}{%
\lemma{trianguli}%
\Bfootnote{%
\textit{(1)}~\textit{6.14.17}  %
\textit{(2)}~\textit{6.4.} %
\textit{(a)}~\textit{19} %
\textit{(b)}~\textit{5}~\textit{L}}}
in quo tres 
\edtext{directiones ponderis\indexS{directiones potentiarum et ponderis} \textit{F},}{%
\lemma{directiones}%
\Bfootnote{%
\textit{(1)}~se 
\textit{(a)}~inters %
\textit{(b)}~intersecant %
\textit{(2)}~ponderis \textit{F},~\textit{L}}}
potentiae \textit{7} et potentiae\indexS{potentia} \textit{8} se intersecant exhibent potentiarum\indexS{potentia} proportiones, 
scilicet unaquaeque potentia exprimitur per latus trianguli cadens in suam directionem,\indexS{directio potentiae} 
seu potentiae \textit{F}, \textit{7}, \textit{8}, sunt ut rectae \textit{6.4}. et \textit{4.5} et \textit{5.6}. Et quia in triangulo latera sunt ut sinus angulorum\indexS{sinus anguli oppositi} 
\edtext{oppositorum, erunt potentiae quaevis ad quamvis,}{%
\lemma{oppositorum,}%
\Bfootnote{%
\textit{(1)}~erunt %
\textit{(2)}~\bars{erit \textit{streicht Hrsg.}} unaquaeque potentia, ut %
\textit{(3)}~erunt potentiae quaevis ad quamvis,~\textit{L}}}
\edtext{ut sinus angulorum\indexS{sinus anguli directionum duarum potentiarum} quos faciunt duae reliquae potentiae.}{%
\lemma{}%
\Afootnote{%
\textit{Zwischen den Zeilen}: Habeat locum, forte licebit concipere \textit{O} et \textit{N} non esse puncta, sed plana, ut coincidant puncta \textit{4}, \textit{5}, \textit{6}.}} 
Et hinc cum regula\indexS{regula generalis pro velocitate progressus} 
\edtext{general\bracks{i}}{\lemma{}\Bfootnote{generalis \textit{L ändert Hrsg.}}}
pro velocitate progressus,\indexS{progressus}\indexS{velocitas progressus} \textit{dm} est velocitas\indexS{velocitas} qua pondus \textit{6} 
\edtext{transit dextrorsum}{%
\lemma{transit}%
\Bfootnote{%
\textit{(1)}~versus %
\textit{(2)}~dextrorsum~\textit{L}}}
vel sinistrorsum, \textit{dy} qua sursum aut 
\edtext{deorsum et erit}{%
\lemma{deorsum,}%
\Bfootnote{%
\textit{(1)}~\textit{dn} est velocitas qua pondus %
\textit{(2)}~et erit~\textit{L}}}
illa ad hanc ut \textit{6.12} ad \textit{8.12}. Si scilicet radius \textit{8.6} moveatur circa centrum 
\edtext{\textit{8},}{%
\lemma{}%
\Bfootnote{%
\textit{8} \textit{erg.~L}}}
velocitas 
\edtext{qua augetur \textit{8.12} sinus complementi\indexS{sinus complementi}}{%
\lemma{qua}%
\Bfootnote{%
\textit{(1)}~transit %
\textit{(2)}~augetur \bars{\textit{8.12} \textit{erg.}} sinus
\textit{(a)}~rectus %
\textit{(b)}~complementi~\textit{L}}}
\textit{m}, (nempe $d\overline{m}$) est ad velocitatem\indexS{velocitas} qua minuitur 
\edtext{\textit{12.6}}{%
\lemma{}%
\Bfootnote{%
\textit{12.6} \textit{erg.~L}}}
\edtext{sinus rectus\indexS{sinus rectus} (seu}{%
\lemma{sinus}%
\Bfootnote{%
\textit{(1)}~comple %
\textit{(2)}~rectus
\bars{\textit{(1)}~$-y$ \textit{(2)}~$-y$ \textit{gestr.}}
\textit{(a)}~(seu ad $-dy$) ut s %
\textit{(b)}~(seu~\textit{L}}}
ad $dy$ nam sinus rectus est $c-y$) ut sinus rectus\indexS{sinus rectus} ad sinum complementi,\indexS{sinus complementi}
\edtext{seu}{%
\lemma{seu}%
\Bfootnote{%
\textit{(1)}~velocitates sunt reciproce ut %
\textit{(2)}~seu~\textit{L}}}
crementa\indexS{crementum} sunt reciproce ut crescentia vel decrescentia. 
Similiter ex velocitatibus\indexS{velocitas} quibus augentur \textit{m} et \textit{y}, 
determinatur velocitas\indexS{velocitas} qua augetur \textit{n} seu $\sqrt{\overline{b-m}^2+\overline{g-y}^2}$.
\edtext{Sumantur $\omega$}{%
\lemma{}\Bfootnote{Sumantur \bars{rectae \textit{gestr.}} $\omega$~\textit{L}}}
et $\beta$ quae sunt respective ad velocitates crescentium\indexS{velocitates crescentium} $m$ et $y$ ut $b-m$, vel $g-y$ ad \textit{n}, summa $\omega + \beta$ erit velocitas\indexS{velocitas decrescendi} qua decrescunt ipsae \textit{n}.
\pend
\pstart
\noindent
\href{\Scan{LH_35_10_14_010r}}%
{\lbrack10~r\textsuperscript{o}\rbrack}\ 
\pend
\begin{edfigure}[0em][1em]
\includegraphics[scale=1]{%
gesamttex/edit_VIII-4/images/LH_35_10_14_009-010_d9_010r.pdf}
\edcaption{\lbrack\textit{Fig.~9}\rbrack}
\end{edfigure}

\pstart
Mirum 
\edtext{autem videri possit,}{%
\lemma{autem}%
\Bfootnote{%
\textit{(1)}~videtur %
\textit{(2)}~videri possit,~\textit{L}}}
quod eadem prodit regula\indexS{regula pro velocitate progressus} ac si coinciderent puncta \textit{4. 5. 6}. 
Hoc generali Methodo ita videtur demonstrari posse, quod magni est momenti. 
Nempe planum aliquod ut \textit{N} eodem modo premitur a pondere\indexS{pondus planum premens} \textit{F}, sive tangat in puncto \textit{O} uno,  
\edtext{sive in linea}{%
\lemma{sive in}%
\Bfootnote{%
\textit{(1)}~pluribus %
\textit{(2)}~linea~\textit{L}}}
recta integra ut \textit{21.22}. Idem est in plano \textit{N}, et ita tantum opus est ex centro gravitatis\indexS{centrum gravitatis} demitti perpendiculares in plana contactus.\indexS{planum contactus} Atque 
\edtext{ita habebitur}{%
\lemma{ita}%
\Bfootnote{%
\textit{(1)}~res proced %
\textit{(2)}~habebitur~\textit{L}}}
exitus et in casu plurium planorum.
\pend
\pstart
Et quae tunc linea directionis centri gravitatis\indexS{linea directionis centri gravitatis}\indexS{centrum gravitatis} non cadit planum verticale\indexS{planum verticale} cum punctis occursuum seu lineis directionum,\indexS{linea directionis}
\edtext{itaque sumemus}{%
\lemma{itaque}%
\Bfootnote{%
\textit{(1)}~duae %
\textit{(2)}~fingemus planum inclinatum composi %
\textit{(3)}~sumemus~\textit{L}}}
planum inclinatum\indexS{planum inclinatum} in quo est centrum gravitatis\indexS{centrum gravitatis} et duae directiones, inde habebitur directio composita\indexS{directio composita} in illo ipso plano inclinato\indexS{planum inclinatum} quae conjuncta cum aliis directionibus\indexS{directio composita} 
\edtext{compositis, unam demum dabit directionem\indexS{directio totalis}}{%
\lemma{compositis,}%
\Bfootnote{%
\textit{(1)}~una demum prodibit directio
\textit{(2)}~unam demum dabit directionem~\textit{L}}}
totalem quae erit ipsa directio centri gravitatis.\indexS{directio centri gravitatis} 
Videndum an non sic demonstrari possit eodem modo premi planum\indexS{corpus planum premens}
\textit{YOZ} sive in puncto \textit{O}, sive in alio quovis a corpore \textit{FON} tangatur. Nempe si tangitur in \textit{F}, ponatur linea rigida \textit{NF3O}, composita ex rectis 
\edtext{\textit{NF3} et \textit{3O}, est}{%
\lemma{\textit{NF3} et}%
\Bfootnote{%
\textit{(1)}~\textit{FO} %
\textit{(2)}~\textit{3O}, %
\textit{(a)}~hujus vectes inflexi ope %
\textit{(b)}~pr %
\textit{(c)}~pondus collectum in \textit{F} premet \bars{planum in \textit{O} \textit{streicht Hrsg.}} certo modo
\textit{(d)}~est~\textit{L}}}
autem \textit{3O} normalis ad 
\edtext{planum \textit{YO}. Apparet}{%
\lemma{planum \textit{YO}.}%
\Bfootnote{%
\textit{(1)}~Patet
\textit{(2)}~Apparet~\textit{L}}}
eodem modo premere, ac fieret vecte inflexo\indexS{vectis inflexus}
\textit{NFZ}, si esset \textit{FZ} normalis ad \textit{OYZ}. 
Et hoc accuratius examinandum\bracks{,} hoc ipso 
\edtext{enim demonstrato, omne}{%
\lemma{enim}%
\Bfootnote{%
\textit{(1)}~demonstratio, omnis Statica %
\textit{(2)}~demonstrato, omne~\textit{L}}}
Staticae\indexS{Statica} negotium\indexS{negotium Staticae} absolutum est. %
\pend%
\begin{edfigure}%
\includegraphics[scale=1]{%
gesamttex/edit_VIII-4/images/LH_35_10_14_009-010_d10_010r.pdf}
\edcaption{\lbrack\textit{Fig.~10}\rbrack}
\end{edfigure}
\pstart
Addendum 
\edtext{jam quid}{%
\lemma{jam}%
\Bfootnote{%
\textit{(1)}~si %
\textit{(2)}~quid fiat, si plures chordae unam in %
\textit{(3)}~quid~\textit{L}}}
fiat si punctum trahatur a pluribus chordis,\indexS{punctum tractum a pluribus chordis} et harum quaevis rursus a pluribus\indexS{chorda trahens tracta a pluribus aliis} quemadmodum et si sustentacula\indexS{sustentaculum} rursus aliis sustentaculis\indexS{sustentaculum alio sustentaculo incumbens} incumbant, sed hac methodo omnia sunt finita.
\pend
\pstart
Res omnis conspirante conclusione etiam inde 
\edtext{conficitur quod punctis corporis cadentibus}{%
\lemma{conficitur}%
\Bfootnote{%
\textit{(1)}~quod corpore tendente linea quacunque, %
\textit{(a)}~qua %
\textit{(b)}~motu aequidistributo 
\textit{(aa)}~seu %
\textit{(bb)}~nempe 
\textit{(cc)}~ita ut %
\textit{(aa)}~per %
\textit{(bb)}~puncta %
\textit{(2)}~quod punctis corporis %
\textit{(a)}~tendentibus ut %
\textit{(b)}~cadentibus~\textit{L}}}
in rectam perpendicularem ad contactum plani in eadem recta tendentibus, posito motu aequidistributo\indexS{motus aequidistributus}
totum corpus eodem modo tendit adeoque et centrum gravitatis.\indexS{centrum gravitatis} 
Et abstrahendo a renisu materiae\indexS{renisus materiae} qui locum habet in vi viva,\indexS{vis viva} nulla oritur gyratio.\indexS{gyratio}
\pend
